\chapter{Órbitas: cuándo y cada cuánto pasa el satélite}\label{cap:orbitas}

\objetivos{
\begin{itemize}
  \item Calcular el periodo de un satélite a partir de su altura.
  \item Explicar qué es una órbita polar heliosíncrona y por qué el satélite cruza cada lugar
  siempre a la misma hora local.
  \item Decir a qué hora observa cada satélite de este libro y qué parte del ciclo diario del
  fuego capta.
  \item Explicar por qué MODIS deja huecos de cobertura cerca del ecuador y VIIRS no.
\end{itemize}}

\section{Por qué importa la órbita}

Un foco de calor es una observación instantánea. Qué fuegos entran en el registro depende, antes
que de cualquier algoritmo, de \emph{cuándo} mira el satélite y de \emph{cuántas veces} mira. Dos
sensores idénticos en órbitas distintas producirían registros distintos. Para comparar MODIS con
VIIRS hay que entender primero cómo se mueven sus plataformas.

\section{Movimiento orbital}

Un satélite en órbita está en caída libre permanente: avanza tan rápido en horizontal que, a
medida que cae, la superficie curva de la Tierra se aleja bajo él al mismo ritmo. Para una
órbita circular, la gravedad proporciona exactamente la fuerza centrípeta necesaria:
\begin{equation}
  \frac{G M m}{r^2} = \frac{m v^2}{r}
  \quad\Longrightarrow\quad
  v = \sqrt{\frac{GM}{r}},
\end{equation}
donde $r$ es la distancia al centro de la Tierra (radio terrestre más altura), y $GM =
\SI{3.986e5}{\kilo\metre\cubed\per\second\squared}$ es el parámetro gravitacional de la Tierra.
La masa del satélite se cancela: todos los objetos a la misma altura orbitan a la misma
velocidad.

El \termino{periodo} $P$, el tiempo de una vuelta completa, es la longitud de la circunferencia
dividida por la velocidad:
\begin{equation}\label{eq:periodo}
  P = \frac{2\pi r}{v} = 2\pi\sqrt{\frac{r^3}{GM}} .
\end{equation}
Esta es la tercera ley de Kepler para órbitas circulares.

\begin{ejemplo}{periodo de los satélites de MODIS y de VIIRS}
Terra y Aqua vuelan a 705~km de altura; Suomi NPP, NOAA-20 y NOAA-21, a unos 824~km. Con el
radio terrestre $R_\oplus = \SI{6378}{\kilo\metre}$, calcula su velocidad, su periodo y cuántas
órbitas dan al día.
\solucion Para MODIS, $r = 6378 + 705 = \SI{7083}{\kilo\metre}$.
\[
 v = \sqrt{\frac{\num{3.986e5}}{7083}} = \SI{7.50}{\kilo\metre\per\second},\qquad
 P = 2\pi\sqrt{\frac{7083^3}{\num{3.986e5}}} = \SI{5933}{\second} = \SI{98.9}{\minute}.
\]
En un día caben $1440/98{,}9 = 14{,}56$ órbitas. Para VIIRS, $r = \SI{7202}{\kilo\metre}$, y se
obtiene $v = \SI{7.44}{\kilo\metre\per\second}$, $P = \SI{101.4}{\minute}$ y 14,20 órbitas al
día. Ambos dan una vuelta al planeta en poco más de hora y media, a unos 27\,000~km/h.
\end{ejemplo}

\section{Órbita polar}

Si el satélite orbitara sobre el ecuador solo vería una franja ecuatorial. Para observar todo el
planeta se usa una \termino{órbita polar}: el satélite pasa cerca de los dos polos en cada
vuelta. La órbita mantiene su orientación en el espacio mientras la Tierra gira debajo, de modo
que en cada vuelta el satélite sobrevuela una franja distinta, desplazada hacia el oeste respecto
de la anterior (figura~\ref{fig:polar}).

\begin{figure}[ht]
\centering
\begin{tikzpicture}
  \begin{scope}
    \shade[ball color=cielo!45] (0,0) circle (2.3);
    \draw[ceniza!70] (-2.3,0) arc (180:360:2.3 and 0.55);
    \draw[ceniza!70,dashed] (2.3,0) arc (0:180:2.3 and 0.55);
    \draw[fuego,very thick] (0.35,2.9) arc (90:-90:0.75 and 2.9);
    \draw[fuego,very thick,dashed] (0.35,-2.9) arc (270:90:0.75 and 2.9);
    \fill[black] (1.02,1.3) circle (0.08);
    \draw[ceniza] (1.12,1.35)--(2.25,1.85);
    \node[anchor=west,font=\footnotesize,fill=white,inner sep=2pt] at (2.25,1.85) {satélite};
    \draw[->,thick] (2.75,0.5) arc (10:-35:1.2) node[midway,right,font=\footnotesize,align=left]{la Tierra\\ gira};
    \draw[->] (0,2.3)--(0,3.3) node[above,font=\footnotesize]{N};
    \node[font=\small\bfseries] at (0,-3.5) {(a) Órbita casi polar};
  \end{scope}
  \begin{scope}[shift={(8.2,0)}]
    \fill[cieloClaro] (-3.8,-2.2) rectangle (3.8,2.2);
    \draw[ceniza] (-3.8,-2.2) rectangle (3.8,2.2);
    \draw[ceniza,dashed] (-3.8,0)--(3.8,0) node[right,font=\footnotesize]{ecuador};
    \foreach \x/\n in {2.6/1,1.3/2,0/3,-1.3/4,-2.6/5}{
      \fill[fuego,opacity=0.22] (\x-0.75,-2.2)--(\x-0.15,-2.2)--(\x+0.55,2.2)--(\x-0.05,2.2)--cycle;
      \draw[fuego,thick,->] (\x-0.45,-2.2)--(\x+0.25,2.2);
      \node[font=\footnotesize] at (\x+0.25,2.45) {órbita \n};
    }
    \draw[<->] (1.2,-1.6)--(-0.1,-1.6) node[midway,below,font=\footnotesize]{$\approx$ 2750 km};
    \node[font=\small\bfseries] at (0,-3.5) {(b) Trazas sucesivas sobre el mapa};
  \end{scope}
\end{tikzpicture}
\caption{Órbita polar. (a) El plano de la órbita permanece casi fijo mientras la Tierra gira
debajo. (b) Vistas en un mapa, las pasadas sucesivas quedan desplazadas hacia el oeste; la franja
sombreada es la zona que el sensor alcanza a ver en cada pasada.}
\label{fig:polar}
\end{figure}

¿Cuánto se desplaza la traza entre una órbita y la siguiente? En el tiempo de una vuelta, la
Tierra gira una fracción $P/\SI{1440}{\minute}$ de su circunferencia. En el ecuador, cuya
longitud es de \SI{40075}{\kilo\metre}:
\begin{equation}\label{eq:separacion}
  d_{\text{ecuador}} = \SI{40075}{\kilo\metre}\times\frac{P}{\SI{1440}{\minute}} .
\end{equation}
Para MODIS resulta $40075 \times 98{,}9/1440 \approx \SI{2750}{\kilo\metre}$; para VIIRS, unos
\SI{2820}{\kilo\metre}. Recuerda estos dos números; volverán en la
sección~\ref{sec:huecos}.

\section{Órbita heliosíncrona: siempre a la misma hora}

Las órbitas de estos satélites no son exactamente polares. Tienen una inclinación de unos
98$^\circ$, es decir, pasan a unos 8$^\circ$ de los polos y giran en sentido ligeramente
contrario a la rotación terrestre. Esa pequeña desviación está elegida con un propósito.

La Tierra no es una esfera perfecta: está achatada, con un abultamiento en el ecuador. Ese
abultamiento tira del plano de la órbita y lo hace girar lentamente, como gira el eje de un
trompo inclinado. El fenómeno se llama \termino{precesión}. Escogiendo bien la altura y la
inclinación, se consigue que el plano orbital dé exactamente una vuelta al año,
$360^\circ/365{,}24 \approx 0{,}986^\circ$ por día, la misma velocidad angular con que la Tierra
gira alrededor del Sol.

\begin{definicion}{órbita heliosíncrona}
Una órbita heliosíncrona es aquella cuyo plano gira al mismo ritmo que la Tierra alrededor del
Sol, de modo que conserva una orientación fija respecto del Sol durante todo el año. Su
consecuencia práctica: el satélite cruza cada latitud siempre a la misma \textbf{hora solar
local}.
\end{definicion}

\begin{figure}[ht]
\centering
\begin{tikzpicture}[scale=0.95]
  \shade[inner color=yellow!80!white,outer color=brasa] (0,0) circle (0.5);
  \node[font=\footnotesize] at (0,-0.8) {Sol};
  \draw[ceniza!50,dashed] (0,0) circle (3.6);
  \foreach \a/\m in {0/marzo,90/junio,180/septiembre,270/diciembre}{
    \begin{scope}[shift={(\a:3.6)}]
      \shade[ball color=cielo!50] (0,0) circle (0.42);
      \draw[fuego,very thick,rotate=\a+90] (-0.85,0)--(0.85,0);
    \end{scope}
  }
  \node[font=\footnotesize,anchor=west] at (4.2,0) {marzo};
  \node[font=\footnotesize,anchor=south] at (0,4.25) {junio};
  \node[font=\footnotesize,anchor=east] at (-4.2,0) {septiembre};
  \node[font=\footnotesize,anchor=north] at (0,-4.25) {diciembre};
  \node[font=\footnotesize,align=left,fuego,anchor=west] at (5.8,1.0) {plano de la órbita\\ (visto de canto)};
  \node[font=\footnotesize,align=left,anchor=west] at (5.8,-0.6) {El plano mantiene el\\ mismo ángulo con la\\ dirección al Sol\\ todo el año.};
\end{tikzpicture}
\caption{Órbita heliosíncrona vista desde encima del polo norte. A lo largo del año el plano
orbital (línea roja) va girando para conservar su orientación respecto del Sol. Por eso el
satélite sobrevuela cada punto a la misma hora local en enero que en julio.}
\label{fig:helio}
\end{figure}

Para el fuego esto es fundamental. Si la hora de paso fuera variando a lo largo del año, el
sensor observaría unas veces el máximo de actividad de la tarde y otras la calma de la mañana, y
sería imposible separar los cambios reales de los cambios de horario. Con una órbita
heliosíncrona, la hora de observación es una constante del sistema.

\begin{definicion}{hora de cruce ecuatorial}
La hora de cruce ecuatorial es la hora solar local a la que el satélite atraviesa el ecuador.
Cada órbita tiene dos cruces, separados doce horas: el \emph{nodo ascendente}, cuando viaja de
sur a norte, y el \emph{nodo descendente}, de norte a sur.
\end{definicion}

\section{A qué hora mira cada satélite}

La tabla~\ref{tab:horas} reúne las horas nominales de cruce.

\begin{table}[ht]
\centering
\begin{tabular}{llccc}
\toprule
Satélite & Sensor & Altura & Cruce diurno & Cruce nocturno \\
\midrule
Terra & MODIS & 705 km & 10:30 & 22:30 \\
Aqua & MODIS & 705 km & 13:30 & 01:30 \\
Suomi NPP & VIIRS & 824 km & 13:30 & 01:30 \\
NOAA-20 & VIIRS & 824 km & 13:30 & 01:30 \\
NOAA-21 & VIIRS & 824 km & 13:30 & 01:30 \\
\bottomrule
\end{tabular}
\caption{Horas solares locales nominales de cruce del ecuador. Los tres satélites VIIRS
comparten la misma órbita, repartidos a lo largo de ella: NOAA-20 va media órbita (unos 50
minutos) separado de Suomi NPP, y NOAA-21 se sitúa entre ambos.}
\label{tab:horas}
\end{table}

Tres observaciones sobre la tabla, todas con consecuencias directas para la armonización.

\textbf{Aqua y los satélites VIIRS son compañeros de horario.} Todos cruzan a las 13:30. Aqua y
Suomi NPP observan los mismos fuegos casi a la misma hora, con diferencias que rara vez superan
unas decenas de minutos. Es la pareja natural para comparar los dos sensores.

\textbf{Terra no tiene sucesor.} Su pasada de las 10:30 capta el fuego de media mañana, y no
existe ningún VIIRS en ese horario. Cuando Terra se apague, esa ventana de observación
desaparecerá.

\textbf{Las horas son nominales.} En sus últimos años de vida, Terra y Aqua han dejado de hacer
las maniobras que mantenían su órbita, y sus horas de cruce se han ido desplazando: la de Terra
hacia más temprano y la de Aqua hacia más tarde. Para los años más recientes de MODIS la hora de
observación ya no es la de la tabla.

\begin{atencion}[La deriva orbital de MODIS]
Desde aproximadamente 2020--2022, Terra y Aqua derivan. Un cambio en la hora de paso altera qué
parte del ciclo diario del fuego se observa, y por tanto el número de focos, sin que el fuego
haya cambiado. Si tu calibración entre MODIS y VIIRS usa los últimos años del solape, parte de
la diferencia que midas se deberá a esta deriva. Es un buen motivo para calibrar con los años
centrales del solape (por ejemplo, 2012--2020) y decirlo expresamente.
\end{atencion}

\section{El ciclo diario del fuego}

¿Por qué tanta insistencia con la hora? Porque el fuego no arde igual a lo largo del día. La
humedad del aire es mínima y la temperatura y el viento son máximos a primera hora de la tarde;
el combustible está más seco y las llamas avanzan con más energía. Además, las quemas agrícolas
se encienden de día. El resultado es un \termino{ciclo diurno} muy marcado, con el máximo entre
las 13 y las 15 horas y el mínimo de madrugada (figura~\ref{fig:diurno}).

\begin{figure}[ht]
\centering
\begin{tikzpicture}
\begin{axis}[width=13.5cm,height=6.6cm,xmin=0,xmax=24,ymin=0,ymax=1.25,
  xlabel={Hora solar local},ylabel={Actividad de fuego (relativa)},
  xtick={0,3,6,9,12,15,18,21,24},ytick={0,0.5,1},yticklabels={0,{0,5},1},
  axis lines*=left,samples=200,domain=0:24]
\addplot[fill=fuego!25,draw=fuego,very thick] {0.04+exp(-((x-14)^2)/(2*2.6^2))} \closedcycle;
\foreach \h/\n/\c/\y in {10.5/Terra/modis/1.12,13.5/{Aqua y VIIRS}/viirs/1.2,22.5/Terra/modis/1.12}{
  \edef\temp{\noexpand\draw[\c,very thick,dashed] (axis cs:\h,0)--(axis cs:\h,\y-0.09);
             \noexpand\node[font=\noexpand\footnotesize,\c] at (axis cs:\h,\y-0.03) {\n};}
  \temp
}
\draw[viirs,very thick,dashed] (axis cs:1.5,0)--(axis cs:1.5,1.11);
\node[font=\footnotesize,viirs,anchor=west] at (axis cs:.25,1.17) {Aqua y VIIRS};
\end{axis}
\end{tikzpicture}
\caption{Forma típica del ciclo diario del fuego de vegetación (curva esquemática) y horas
nominales de paso. La pasada de las 13:30 cae cerca del máximo; la de las 10:30 encuentra los
fuegos todavía creciendo; las nocturnas ven solo lo que sobrevive a la noche.}
\label{fig:diurno}
\end{figure}

La forma exacta de la curva cambia con la región y con el tipo de fuego. Las quemas de residuos
agrícolas producen un pico estrecho, casi sin actividad nocturna; los grandes incendios
forestales tienen un ciclo más ancho y siguen ardiendo de noche.

\begin{ejemplo}{lo que ve cada pasada}
Supón que en una región la actividad sigue la curva de la figura~\ref{fig:diurno}, que
aproximaremos por una campana centrada en las 14:00 con desviación de 2,6 horas. ¿Qué fracción
del máximo observan la pasada de las 10:30 y la de las 13:30?
\solucion La campana vale $\exp[-(h-14)^2/(2\cdot 2{,}6^2)]$. A las 13:30, $h - 14 = -0{,}5$ y
el valor es $\exp(-0{,}018) = 0{,}98$. A las 10:30, $h-14 = -3{,}5$ y el valor es
$\exp(-0{,}906) = 0{,}40$. La pasada de la mañana encuentra menos de la mitad de la actividad
que la de la tarde. Un sensor que solo pasara a las 10:30 contaría, en esta región, alrededor
del 40\,{}\% de lo que cuenta uno que pasa a las 13:30, aunque ambos fueran idénticos.
\end{ejemplo}

\section{Franja de observación y huecos de cobertura}\label{sec:huecos}

El sensor no observa solo el punto que tiene debajo. Barre una franja perpendicular a su
trayectoria, la \termino{franja de barrido} o \emph{swath}, cuyo ancho depende de la altura y
del ángulo máximo con que mira hacia los lados. El punto situado exactamente bajo el satélite se
llama \termino{nadir}.

\begin{center}
\begin{tabular}{lcc}
\toprule
 & MODIS & VIIRS \\
\midrule
Ancho de la franja & 2330 km & 3040 km \\
Separación entre órbitas en el ecuador & 2750 km & 2820 km \\
\bottomrule
\end{tabular}
\end{center}

Compara las dos filas. La franja de MODIS es \emph{más estrecha} que la distancia entre trazas
consecutivas: en el ecuador, entre una pasada y la siguiente queda un hueco de unos 420~km que
ese día no se observa (figura~\ref{fig:huecos}). La franja de VIIRS es \emph{más ancha} que la
separación: las pasadas se solapan y no queda ningún hueco.

\begin{figure}[ht]
\centering
\begin{tikzpicture}[x=0.00175cm,y=1cm]
  \begin{scope}
    \node[font=\small\bfseries,modis] at (2750,2.1) {MODIS};
    \foreach \c in {0,2750,5500}{
      \fill[modis,opacity=0.3] (\c-1165,0) rectangle (\c+1165,1.4);
      \draw[modis,thick] (\c,0)--(\c,1.4);
    }
    \foreach \c in {1375,4125}{
      \fill[pattern=north east lines,pattern color=fuego] (\c-210,0) rectangle (\c+210,1.4);
    }
    \draw[<->] (0,-0.3)--(2750,-0.3) node[midway,below,font=\footnotesize]{2750 km};
    \draw[<->] (-1165,1.65)--(1165,1.65) node[midway,above,font=\footnotesize]{2330 km};
    \node[font=\footnotesize,fuego] at (4125,-0.35) {hueco};
  \end{scope}
  \begin{scope}[shift={(0,-4.2)}]
    \node[font=\small\bfseries,viirs] at (2820,2.1) {VIIRS};
    \foreach \c in {0,2820,5640}{
      \fill[viirs,opacity=0.25] (\c-1520,0) rectangle (\c+1520,1.4);
      \draw[viirs,thick] (\c,0)--(\c,1.4);
    }
    \draw[<->] (0,-0.3)--(2820,-0.3) node[midway,below,font=\footnotesize]{2820 km};
    \draw[<->] (-1520,1.65)--(1520,1.65) node[midway,above,font=\footnotesize]{3040 km};
    \node[font=\footnotesize] at (4230,0.7) {solape};
  \end{scope}
\end{tikzpicture}
\caption{Cobertura en el ecuador de tres órbitas consecutivas, a escala. MODIS deja franjas sin
observar (rayadas) entre pasadas; las franjas de VIIRS se solapan.}
\label{fig:huecos}
\end{figure}

A medida que aumenta la latitud las trazas convergen hacia los polos y los huecos de MODIS se
cierran; por encima de unos 30$^\circ$ ya no existen. Pero el Perú está entre 0$^\circ$ y
18$^\circ$ de latitud sur: de lleno en la zona de huecos.

\begin{intuicion}[Qué significa un hueco]
Los huecos no caen siempre en el mismo sitio; se desplazan de un día a otro. Para un lugar dado
cerca del ecuador, lo que ocurre es que algunos días, aproximadamente uno de cada seis o siete,
Aqua no lo observa de día. Ese día cualquier fuego en ese lugar es invisible para Aqua, pero no
para VIIRS. La figura~\ref{fig:worldview} del capítulo~\ref{cap:modisviirs} muestra un caso
real sobre el centro del Perú.
\end{intuicion}

\section{Número de observaciones por día}

Juntando todo, para un lugar de latitud baja como Junín, el número de oportunidades de
observación por día ha cambiado mucho a lo largo del registro:

\begin{center}
\begin{tabular}{lcl}
\toprule
Periodo & Pasadas diarias (día + noche) & Sensores \\
\midrule
2000--2002 & hasta 2 & Terra \\
2002--2011 & hasta 4 & Terra, Aqua \\
2012--2017 & hasta 4 MODIS + 2 VIIRS & + Suomi NPP \\
2018--2023 & hasta 4 MODIS + 4 VIIRS & + NOAA-20 \\
2024-- & hasta 4 MODIS + 6 VIIRS & + NOAA-21 \\
\bottomrule
\end{tabular}
\end{center}

<<Hasta>> porque los huecos de MODIS reducen el número real. La consecuencia es directa: si se
suman los focos de todos los satélites disponibles en cada época, el total crece simplemente
porque hay más ojos mirando.

\begin{parareto}
Antes de calibrar nada, \textbf{fija el conjunto de satélites} que entra en la serie. Una opción
limpia: usar solo Aqua para la era MODIS y solo Suomi NPP para la era VIIRS, que comparten la
hora de las 13:30. La columna \texttt{satellite} de los datos permite filtrar. Añadir NOAA-20 y
NOAA-21 mejora la detección reciente, pero rompe la comparabilidad con los años anteriores a
menos que se corrija expresamente.
\end{parareto}

\resumen

\begin{itemize}
  \item El periodo de una órbita circular depende solo de su altura:
  $P = 2\pi\sqrt{r^3/GM}$. MODIS orbita en unos 99 minutos y VIIRS en unos 101.
  \item Una órbita \textbf{polar} permite ver todo el planeta; una \textbf{heliosíncrona}
  garantiza que el satélite pase siempre a la misma hora solar local.
  \item \textbf{Terra} cruza a las 10:30/22:30; \textbf{Aqua} y los tres \textbf{VIIRS}, a las
  13:30/01:30. Aqua y Suomi NPP forman la pareja natural de comparación. Terra no tiene sucesor.
  \item El fuego tiene un \textbf{ciclo diurno} con máximo a primera hora de la tarde. La hora
  de paso determina qué fracción de la actividad se observa.
  \item Las órbitas de Terra y Aqua \textbf{derivan} en sus últimos años.
  \item La franja de MODIS (2330~km) deja \textbf{huecos} en latitudes bajas; la de VIIRS
  (3040~km) no.
  \item El número de satélites en vuelo ha cambiado; hay que fijar cuáles se usan.
\end{itemize}

\ejercicios

\begin{enumerate}
  \item Calcula el periodo de un satélite a 500~km de altura y el de uno a 1000~km. ¿Cuántas
  órbitas diarias da cada uno?
  \item Con la ecuación~\eqref{eq:separacion}, calcula la separación entre trazas consecutivas
  en el ecuador para un satélite con periodo de 95 minutos. ¿Qué ancho mínimo de franja necesita
  para no dejar huecos?
  \item La separación entre trazas a una latitud $\varphi$ es aproximadamente la del ecuador
  multiplicada por $\cos\varphi$. ¿A partir de qué latitud dejan de existir los huecos de
  MODIS? ¿Y a la latitud de Huancayo ($12^\circ$~S), cuánto mide el hueco?
  \item Con la campana del ejemplo de este capítulo, calcula qué fracción del máximo observa una
  pasada a las 15:00 y otra a las 09:00. Si Aqua deriva de las 13:30 hacia las 15:00, ¿sube o
  baja su conteo en esa región? ¿Y si Terra deriva de las 10:30 hacia las 09:00?
  \item Un compañero propone sumar los focos de Terra y Aqua para la era MODIS y los de Suomi
  NPP, NOAA-20 y NOAA-21 para la era VIIRS, <<para tener el máximo de datos>>. Escribe dos
  objeciones.
  \item ¿Por qué una órbita heliosíncrona necesita una inclinación mayor de 90$^\circ$ y no
  exactamente 90$^\circ$? Explícalo en un párrafo usando la idea de precesión.
\end{enumerate}
